%% ── Section 5: Results and Discussion ───────────────────────
\section{Results and Discussion}
\label{sec:results}

\subsection{Peak Diagnostic Performance}

Table~\ref{tab:results} reports the peak validation QWK achieved over 10
communication rounds. Under the ideal \textbf{IID} setting, all federated algorithms
reach comparable peak QWK values (0.8357--0.8386), remaining within 0.014 of the
centralised ceiling (0.8494), confirming that algorithmic overhead is negligible
when data is evenly distributed.

On the clinically modeled \textbf{Real-Site} partition, OrdinalFed achieves the
highest peak QWK of \textbf{0.8325}, outperforming FedProx (0.8312) and FedAvg
(0.8223), closing the gap to the centralised ceiling to just 1.69~pp. Under extreme
\textbf{Dirichlet ($\alpha=0.1$)} skew, all three methods achieve near-identical
peak QWK ($\approx$0.799), owing to the strong warm-start initialisation. However,
peak performance is misleading in Non-IID regimes; the fundamental advantage of
OrdinalFed lies in its resilience against catastrophic forgetting.

\begin{table}[t]
\caption{Peak Validation QWK Across 10 Communication Rounds.
Bold: best FL method per column.}
\label{tab:results}
\centering
\setlength{\tabcolsep}{5pt}
\begin{tabular}{@{}lccc@{}}
\toprule
\textbf{Method} &
\textbf{IID} &
\textbf{Dirichlet ($\alpha=0.1$)} &
\textbf{Real-Site} \\
\midrule
Centralised (Ceiling) & 0.8494 & 0.8494 & 0.8494 \\
\midrule
FedAvg~\cite{mcmahan2017fedavg}
  & \textbf{0.8386} & \textbf{0.7993} & 0.8223 \\
FedProx~\cite{li2020fedprox} ($\mu=0.01$)
  & 0.8357 & 0.7987 & 0.8312 \\
\textbf{OrdinalFed (Ours)}
  & 0.8372 & 0.7988 & \textbf{0.8325} \\
\bottomrule
\end{tabular}
\end{table}

\subsection{Resilience to Catastrophic Forgetting}

As visualised in the central panel of Fig.~\ref{fig:qwk_curves}, FedAvg suffers
catastrophic collapse under extreme skew ($\alpha=0.1$), plunging from 0.7870 to
0.4064 in round~9. This occurs because Client~0's monochromatic Grade~0 dataset
forces its local gradients to predict all inputs as Grade~0. When aggregated via
sample counts, this corrupted update shatters the global decision boundaries.
FedProx dampens this instability partially by tethering local weights to the global
model but still crashes to 0.5963 in round~10. In contrast, OrdinalFed maintains
consistent diagnostic performance, never falling below a QWK floor of \textbf{0.6238}.

As quantified in Table~\ref{tab:stability}, OrdinalFed achieves the narrowest
peak-to-trough fluctuation range (\textbf{0.1751} vs.\ 0.2024 for FedProx and
0.3929 for FedAvg), ensuring the round-to-round consistency required for reliable
clinical deployment.

\begin{table}[t]
\caption{Stability Metrics Under Extreme Dirichlet Skew ($\alpha=0.1$).
Bold: best value per column ($\uparrow$ Peak/Worst/Mean; $\downarrow$ Range).}
\label{tab:stability}
\centering
\setlength{\tabcolsep}{4pt}
\begin{tabular}{@{}lcccc@{}}
\toprule
\textbf{Method} &
\textbf{Peak} &
\textbf{Worst} &
\textbf{Mean} &
\textbf{Range} \\
\midrule
FedAvg~\cite{mcmahan2017fedavg}
  & \textbf{0.7993} & 0.4064 & 0.7286 & 0.3929 \\
FedProx~\cite{li2020fedprox}
  & 0.7987 & 0.5963 & \textbf{0.7524} & 0.2024 \\
\textbf{OrdinalFed (Ours)}
  & 0.7988 & \textbf{0.6238} & 0.7453 & \textbf{0.1751} \\
\bottomrule
\end{tabular}
\end{table}

\subsection{Mechanistic Analysis of Stability}

OrdinalFed's stability is driven by two explicit mechanisms acting jointly.
First, by penalising squared grade distance, \textbf{OrdinalLoss} prevents local
representations from collapsing into single-class degenerate states. The severe
penalty for Grade~4 $\to$ Grade~0 misclassifications ensures the network retains
its ordinal hierarchy even when trained exclusively on Grade~0 images. Second,
under $\alpha=0.1$, Client~0 produces severe validation degradation on $\mathcal{V}$
($\kappa_0^t \to 0$). As $\beta(t) \to 0$ in later rounds,
\textbf{QWK-Weighted Aggregation} automatically suppresses Client~0's weight
$\omega_0^t$ to near zero, dynamically quarantining the corrupted update before it
can poison the global model --- the mechanism that prevents the collapse seen in
FedAvg and FedProx.

\subsection{Clinical Relevance in Real-Site Deployments}

In the real-site setting (Fig.~\ref{fig:qwk_curves}, right panel), OrdinalFed
delivers a sustained 1.02~pp gain over FedAvg, yielding a 10-round mean QWK of
0.8261 versus 0.8176 for FedAvg and 0.8241 for FedProx. The specialist referral
centre (Client~2, G3/G4-heavy, 643 images) receives moderated aggregation weight
due to sub-optimal balanced validation performance ($\kappa_2^t \in [0.57, 0.74]$),
whereas the balanced urban hospital (Client~4) is up-weighted despite its smaller
sample size (9\%). This demonstrates clinically rational aggregation that aligns
model influence with diagnostic quality rather than raw data volume.

\subsection{Communication and Computational Overhead}

Because OrdinalFed evaluates QWK scores strictly on the server using proxy set
$\mathcal{V}$, it requires zero additional communication overhead compared to
standard FedAvg. The computational burden is limited to exactly $K \times 50$
forward passes on the server per round, executing in under 200~ms on standard
hardware, making OrdinalFed highly scalable to larger federations.

\subsection{Limitations}

Due to severe natural dataset scarcity in the DDR benchmark (Grade~1 represents
only 4.9\% of the total training volume), Grade~1 sensitivity remains below 0.15
across all tested methods. Furthermore, the 50-image server set $\mathcal{V}$
introduces minor QWK estimation noise in final rounds when $\beta(t) \to 0$.
Setting a non-zero $\beta$ lower bound and expanding $\mathcal{V}$ via generative
augmentation are promising directions for future work.
